% Part: first-order-logic
% Chapter: completeness
% Section: henkin-expansion

\documentclass[../../../include/open-logic-section]{subfiles}

\begin{document}

\olfileid[pt]{fol}{com}{hen}

\olsection{Expansão de Henkin}

\begin{explain}
Parte do desafio em provar o teorema da completude é que o
modelo que construímos a partir de um conjunto consistente completo~$\Gamma$ deve tornar
verdadeiras todas as !!{formula}s quantificadas em~$\Gamma$.  Para
garantir isso, usamos um truque devido a Leon Henkin.  Em essência, o
truque consiste em expandir a linguagem por infinitos !!{constant}s
e adicionar, para cada !!{formula} com uma !!{variable} livre $!A(x)$, uma
fórmula da forma
\iftag{prvEx}
      {$\lexists[x][!A(x)] \lif !A(c)$}
      {$\lnot\lforall[x][!A(x)] \lif \lnot !A(c)$},
onde $c$ é um dos novos !!{constant}s.  Quando construímos a
!!{structure} que satisfaz~$\Gamma$, isso garantirá que cada
\iftag{prvEx}
{sentença existencial verdadeira tenha uma testemunha}
{sentença universal falsa tenha um contraexemplo}
entre as novas constantes.
\end{explain}

\begin{prop}
\ollabel{prop:lang-exp}
Se $\Gamma$ é consistente em $\Lang L$ e $\Lang L'$ é obtida de
$\Lang L$ adicionando-se !!a{denumerable} conjunto de novos !!{constant}s $\Obj d_0$,
$\Obj d_1$, \dots, então $\Gamma$ é consistente em~$\Lang L'$.
\end{prop}

\begin{defn}[Conjunto saturado]
Um conjunto $\Gamma$ de !!{formula}s de uma linguagem $\Lang {L}$ é
\emph{saturado} se e somente se, para cada !!{formula}~$!A(x) \in \Frm[L]$ com uma
!!{variable} livre~$x$, existe !!a{constant}~$c \in \Lang{L}$ tal
que
\iftag{prvEx}
      {$\lexists[x][!A(x)] \lif !A(c) \in \Gamma$}
      {$\lnot\lforall[x][!A(x)] \lif \lnot !A(c) \in \Gamma$}.
\end{defn}

A definição a seguir será usada na prova do próximo teorema.

\begin{defn}
\ollabel{defn:henkin-exp}
Seja $\Lang L'$ como em \olref{prop:lang-exp}.  Fixe uma enumeração
$!A_0(x_0)$, $!A_1(x_1)$, \dots de todas as !!{formula}s~$!A_i(x_i)$
de~$\Lang L'$ nas quais uma variável ($x_i$) ocorre livre.  Definimos
as !!{sentence}s~$!D_n$ por indução sobre~$n$.

Seja $c_0$ o primeiro !!{constant} dentre os $\Obj d_i$ que adicionamos
a~$\Lang{L}$ que não ocorre em~$!A_0(x_0)$.  Supondo que
$!D_0$, \dots,~$!D_{n-1}$ já tenham sido definidas, seja $c_n$ o
primeiro dentre os novos !!{constant}s~$\Obj d_i$ que não ocorre nem em
$!D_0$, \dots,~$!D_{n-1}$ nem em~$!A_n(x_n)$.

Agora seja $!D_{n}$ a !!{formula} \iftag{prvEx}
{$\lexists[x_{n}][!A_{n}(x_{n})] \lif
  !A_{n}(c_{n})$}{$\lnot\lforall[x_{n}][!A_{n}(x_{n})] \lif \lnot
  !A_{n}(c_{n})$}.
\end{defn}

\begin{lem}
\ollabel{lem:henkin}
Todo conjunto consistente~$\Gamma$ pode ser estendido a um conjunto saturado
consistente~$\Gamma'$.
\end{lem}

\begin{proof}
Dado um conjunto consistente de sentenças~$\Gamma$ em uma linguagem~$\Lang{L}$,
expanda a linguagem adicionando !!a{denumerable} conjunto de novos
!!{constant}s para formar~$\Lang{L'}$.  Por \olref{prop:lang-exp}, $\Gamma$
ainda é consistente na linguagem mais rica.  Além disso, seja $!D_i$ como
em \olref{defn:henkin-exp}.  Seja
\begin{align*}
\Gamma_0 & = \Gamma \\
\Gamma_{n+1} & = \Gamma_n \cup \{!D_n \}
\end{align*}
isto é, $\Gamma_{n+1} = \Gamma \cup \{ !D_0, \dots, !D_n \}$, e seja
$\Gamma' = \bigcup_{n} \Gamma_n$.  $\Gamma'$ é claramente saturado.

Se $\Gamma'$ fosse inconsistente, então, para algum $n$, $\Gamma_n$ seria
inconsistente (Exercício: explique por quê).  Portanto, para mostrar que $\Gamma'$ é
consistente, basta mostrar, por indução sobre~$n$, que cada
conjunto~$\Gamma_n$ é consistente.

A base da indução é simplesmente a afirmação de que $\Gamma_0 = \Gamma$ é
consistente, que é a hipótese do teorema.  Para o passo de indução,
suponha que $\Gamma_{n}$ seja consistente mas $\Gamma_{n+1} =
\Gamma_n \cup \{!D_n\}$ seja inconsistente.  Lembre que $!D_n$~é
\iftag{prvEx}
{$\lexists[x_{n}][!A_{n}(x_n)] \lif !A_{n}(c_{n})$}
{$\lnot\lforall[x_{n}][!A_{n}(x_n)] \lif \lnot !A_{n}(c_{n})$},
onde $!A_n(x_n)$ é !!a{formula} de $\Lang{L'}$ com apenas a
variável~$x_n$ livre. Pela maneira como escolhemos os~$c_n$ (ver
\olref{defn:henkin-exp}), $c_n$ não ocorre em~$!A_n(x_n)$ nem
em~$\Gamma_n$.

Se $\Gamma_n \cup \{!D_n\}$ é inconsistente, então $\Gamma_n
\Proves \lnot !D_n$ e, portanto, ambos os seguintes valem:
\iftag{prvEx}{
\[
\Gamma_n \Proves \lexists[x_n][!A_n(x_n)]
\qquad
\Gamma_n \Proves \lnot !A_n(c_n)
\]}{
\[
\Gamma_n \Proves \lnot\lforall[x_n][!A_n(x_n)]
\qquad
\Gamma_n \Proves !A_n(c_n)
\]}
Como $c_n$ não ocorre em
$\Gamma_n$ ou em~$!A_n(x_n)$, 
\tagrefs{prfAX/{fol:axd:qpr:thm:strong-generalization},prfSC/{fol:seq:qpr:thm:strong-generalization},prfND/{fol:ntd:qpr:thm:strong-generalization},prfTab/{fol:tab:qpr:thm:strong-generalization}} se aplica.
De \iftag{prvEx}
{$\Gamma_n \Proves \lnot !A_n(c_n)$}
{$\Gamma_n \Proves !A_n(c_n)$},
obtemos
\iftag{prvEx}
{$\Gamma_n \Proves \lforall[x_n][\lnot !A_n(x_n)]$}
{$\Gamma_n \Proves \lforall[x_n][!A_n(x_n)]$}.
Assim, temos que ambos
\iftag{prvEx}
{$\Gamma_n \Proves \lexists[x_n][!A_n(x_n)]$ e
$\Gamma_n \Proves \lforall[x_n][\lnot !A_n(x_n)]$}
{$\Gamma_n \Proves \lnot\lforall[x_n][!A_n(x_n)]$ e
$\Gamma_n \Proves \lforall[x_n][!A_n(x_n)]$},
de modo que o próprio $\Gamma_n$ é inconsistente.
\iftag{prvEx}
{(Note que
\iftag{prvAll}
{$\lforall[x_n][\lnot !A_n(x_n)] \Proves
  \lnot\lexists[x_n][!A_n(x_n)]$}
{$\lforall[x_n][\lnot !A_n]$ é definida como
  $\lnot\lexists[x_n][\lnot\lnot !A_n(x_n)]$ e
  $\lnot\lexists[x_n][\lnot\lnot !A_n(x_n)] \Proves
  \lnot\lexists[x_n][!A_n(x_n)]$}.)}{}
Contradição: supunha-se que $\Gamma_n$ fosse consistente.  Logo,
$\Gamma_n \cup \{ !D_n\}$ é consistente.
\end{proof}

\begin{explain}
Mostraremos agora que conjuntos \emph{completos}, consistentes que são
saturados têm a propriedade de que \iftag{prvAll}{contêm uma
  !!{sentence} universalmente quantificada se e somente se contêm todas as suas
  instâncias}{}\iftag{defAll,defEx}{}{ e }\iftag{prvAll}{contêm uma
  !!{sentence} existencialmente quantificada se e somente se contêm pelo menos uma
  instância}{}. Usaremos isso para mostrar que a !!{structure} que
geraremos a partir de um conjunto completo, consistente, saturado torna todas as suas
sentenças quantificadas verdadeiras.
\end{explain}

\begin{prop}\ollabel{prop:saturated-instances}
Suponha que $\Gamma$ seja completo, consistente e saturado.
\begin{tagenumerate}{prvEx,prvAll}
\tagitem{prvEx}{$\lexists[x][!A(x)] \in \Gamma$ se e somente se $!A(t) \in \Gamma$
  para pelo menos um termo fechado~$t$.}{}
\tagitem{prvAll}{$\lforall[x][!A(x)] \in \Gamma$ se e somente se $!A(t) \in \Gamma$
  para todo termo fechado~$t$.}{}
\end{tagenumerate}
\end{prop}

\begin{proof}
  \begin{tagenumerate}{prvEx,prvAll}
  \tagitem{prvEx}{%
    \iftag{probEx}{Exercício.}{Suponha primeiro que $\lexists[x][!A(x)]
      \in \Gamma$.  Como $\Gamma$ é saturado,
      $(\lexists[x][!A(x)] \lif !A(c)) \in \Gamma$ para algum
      !!{constant}~$c$. Por
      \tagrefs{prfAX/{fol:axd:ppr:prop:provability-lif},prfSC/{fol:seq:ppr:prop:provability-lif},prfND/{fol:ntd:ppr:prop:provability-lif},prfTab/{fol:tab:ppr:prop:provability-lif}}, item~(1),
      e \olref[ccs]{prop:ccs}\olref[ccs]{prop:ccs-prov-in}, $!A(c)
      \in \Gamma$.

    Para a outra direção, a saturação não é necessária: Suponha
    $!A(t) \in \Gamma$. Então $\Gamma \Proves \lexists[x][!A(x)]$ por
      \tagrefs{prfAX/{fol:axd:qpr:prop:provability-quantifiers},prfSC/{fol:seq:qpr:prop:provability-quantifiers},prfND/{fol:ntd:qpr:prop:provability-quantifiers},prfTab/{fol:tab:qpr:prop:provability-quantifiers}}, item~(1). Por
    \olref[ccs]{prop:ccs}\olref[ccs]{prop:ccs-prov-in},
    $\lexists[x][!A(x)] \in \Gamma$.}}{}
  \tagitem{prvAll}{%
    \iftag{probAll}{Exercício.}{Suponha que $!A(t) \in \Gamma$ para
      todo termo fechado~$t$. Por contradição, assuma
      $\lforall[x][!A(x)] \notin \Gamma$. Como $\Gamma$ é completo,
      $\lnot\lforall[x][!A(x)] \in \Gamma$. Por saturação,
      \iftag{prvEx}{$(\lexists[x][\lnot !A(x)] \lif \lnot !A(c)) \in
        \Gamma$}{$(\lnot\lforall[x][!A(x)] \lif \lnot !A(c)) \in
        \Gamma$} para algum !!{constant}~$c$. Por hipótese, como $c$
      é um termo fechado, $!A(c) \in \Gamma$. Mas isso tornaria
      $\Gamma$ inconsistente. (Exercício: dê a !!{derivation} que mostra que
      \[
      \lnot \lforall[x][!A(x)],
      \iftag{prvEx}{\lexists[x][\lnot !A(x)] \lif \lnot
        !A(c)}{\lnot\lforall[x][!A(x)] \lif \lnot
        !A(c)}, !A(c)
      \]
      é inconsistente.)

      Para a direção reversa, não precisamos de saturação: Suponha
      $\lforall[x][!A(x)] \in \Gamma$.  Então $\Gamma \Proves !A(t)$
      por
      \tagrefs{prfAX/{fol:axd:qpr:prop:provability-quantifiers},prfSC/{fol:seq:qpr:prop:provability-quantifiers},prfND/{fol:ntd:qpr:prop:provability-quantifiers},prfTab/{fol:tab:qpr:prop:provability-quantifiers}},
      item~(2). Obtemos $!A(t) \in \Gamma$ por
      \olref[ccs]{prop:ccs}.}}{}
  \end{tagenumerate}
\end{proof}

\end{document}
